
\documentclass{beamer}
\usepackage{beamerthemeshadow}
\usepackage[latin1]{inputenc}
\usepackage[T1]{fontenc}
\usepackage{etex}
\usepackage[all,knot,arc,frame,matrix,arrow,graph,curve,color]{xy}
\usepackage{amssymb}
\usepackage{leftidx}
\usepackage{graphics}
\usepackage{algorithmic}
\usepackage{algorithm}



\newxyColor{pink}{1.0 0.4 0.5}{rgb}{}

\begin{document}


\title{Teoria de Grafos para Computa��o}  
\subtitle{(Algoritmos\footnote{Alguns algoritmos baseados em http://www.ic.uff.br/\~celso/})}
\author{Prof. Silvio Jamil F. Guimar�es}
\date{2010-2} 

\frame{\titlepage} 


%%%%%%%%%%%%%%%%%%%%%%%%%%%%%%%%%%%%%%%%%%%%%%%%%%%%%%%%%%%%%%%%%%%%%%%%%%%%%%%%%%%%
%%%==================================SLIDE=========================================
%%%%%%%%%%%%%%%%%%%%%%%%%%%%%%%%%%%%%%%%%%%%%%%%%%%%%%%%%%%%%%%%%%%%%%%%%%%%%%%%%%%%

%\frame{\frametitle{Table of contents}\tableofcontents} 

%\section{Matriz de incid�ncia}
\begin{frame}{Caminhamento em grafos}
\begin{block}{Busca em profundidade}
Caminha no grafo visitando todos os seus v�rtices sempre procurando o v�rtice mais profundo.
\end{block}
\begin{block}{Busca em largura}
Expandir o conjunto de v�rtices de forma uniforme em que s�o visitados todos os v�rtices de mesma dist�ncia ao in�cio antes de visitar outros n�veis.
\end{block}


\only<1>{
\begin{block}{}
\[\xygraph{
!{<0cm,0cm>;<1.0cm,0cm>:<0cm,1.0cm>::}
!{( 0, 0)}*+[o][F-]{c}="c"
!{( 0, 1)}*+[o][F-]{a}="a"
!{(1.0,1)}*+[o][F-]{e}="e"
!{(1,-0.5)}*+[o][F-]{d}="d"
!{(-1,-0.5)}*+[o][F-]{b}="b"
  "a":"b" "a":"c" "a":"d"
  "c":"b" "c":"d" "e":"d"
}\]
\end{block}
}


\end{frame}











\end{document}